\documentclass[../../../uem_mat_mei_q.tex]{subfiles}

\begin{document}
    \ref{q:uem_mat_mei_2025_icom_03_q:1}\,～\,\ref{q:uem_mat_mei_2025_icom_03_q:3}\,%
    の空欄 $\myboxd{テ}$ ～ $\myboxd{ニ}$ については下の解答群から解答を選び，%
    所定の解答欄（表面）にその番号をマークせよ．\ref{q:uem_mat_mei_2025_icom_03_q:4}\,%
    については所定の解答欄（裏面）に解答経過と答をともに記せ．

    \vspace{7pt}
    曲線$C:y=x^2+2x+a^2$とする．ただし$a$は実数で$a>1$とする．また，%
    原点と曲線$C$の頂点を通る直線を$l_1$，曲線$C$の接線のうち，原点を通りかつ傾きが正である%
    ものを$l_2$とする．このとき，次の問に答えよ．

    \begin{enumerate}
        \item%
            曲線$C$の頂点の座標を$a$を用いて表すと，\\
            $\left(\,\myboxd{テ}\,,\;\myboxd{ト}\,\right)$である．
            \label{q:uem_mat_mei_2025_icom_03_q:1}
        \item%
            直線$l_1$の式を$a$を用いて表すと，$y=\left(\,\myboxd{ナ}\,\right)x$\\
            である．
        \item%
            曲線$C$と直線$l_2$の接点の$x$座標を$a$を用いて表すと，\\
            $\myboxd{ニ}$である．
            \label{q:uem_mat_mei_2025_icom_03_q:3}
        \item%
            曲線$C$と直線$l_1$で囲まれた部分の面積を$S_1$，曲線$C$と直線$l_1$および$l_2$%
            で囲まれた部分の面積を$S_2$とするとき，$S_1=S_2$となる$a$の値を求めよ．
            \label{q:uem_mat_mei_2025_icom_03_q:4}
    \end{enumerate}

    \vspace{6pt}
    $\myboxd{テ}$ ～ $\myboxd{ニ}$の解答群
    
    （同じものを繰り返し選んでもよい）

    \begin{enumerate}[topsep=0pt, itemsep=0pt, itemindent=0pt, leftmargin=1\zw, labelsep=0pt]
        \item[　]%
            \setlength{\fboxrule}{1pt}\setlength{\fboxsep}{1pt}%
            \fbox{\setlength{\fboxrule}{.5pt}\setlength{\fboxsep}{5pt}%
            \fbox{
                \begin{minipage}{\dimexpr\linewidth-2\zw}%
                    　\makebox[\dimexpr \linewidth/3-\zw/3][l]{\myegga{0}\myspacea$-1$}%
                    \makebox[\dimexpr \linewidth/3-\zw/3][l]{\myegga{1}\myspacea$1$}%
                    \makebox[\dimexpr \linewidth/3-\zw/3][l]{\myegga{2}\myspacea$-a$}

                    　\makebox[\dimexpr \linewidth/3-\zw/3][l]{\myegga{3}\myspacea$a$}%
                    \makebox[\dimexpr \linewidth/3-\zw/3][l]{\myegga{4}\myspacea$-a-1$}%
                    \makebox[\dimexpr \linewidth/3-\zw/3][l]{\myegga{5}\myspacea$a-1$}
                    
                    　\makebox[\dimexpr \linewidth/3-\zw/3][l]{\myegga{6}\myspacea$-a^2-1$}%
                    \makebox[\dimexpr \linewidth/3-\zw/3][l]{\myegga{7}\myspacea$-a^2+1$}%
                    \makebox[\dimexpr \linewidth/3-\zw/3][l]{\myegga{8}\myspacea$a^2-1$}
                    
                    　\myegga{9}\myspacea$a^2+1$
                \end{minipage}%
            }}
    \end{enumerate}
\end{document}